%------------------------------------------------------------------------------%
\begin{question}
  Esboce a trajetória definida pela curva paramétrica
  \[
  \begin{cases}
    x = t \cos(t\pi) \\
    y = t \sin(t\pi)
  \end{cases}
  \]
  \(0\leq t\)
\end{question}

%------------------------------------------------------------------------------%
\begin{resolution}{Solução}
  Similar a parametrização do circulo trigonométrico

  \pause
  \begin{ceqn}
  \[
  \renewcommand*{\arraystretch}{1.7}
  \begin{array}{c<{\quad}c<{\quad}c<{\quad}c<{\quad}c<{\quad}c}
    \hline
    t            & t\pi            & \cos(t\pi) & \sin(t\pi) & x  & y             \\ \hline
    0            & 0               & 1          & 0          & 0  & 0             \\
    \dfrac{1}{2} & \dfrac{\pi}{2}  & 0          & 1          & 0  & \dfrac{1}{2}  \\
    1            & \pi             & -1         & 0          & -1 & 0             \\
    \dfrac{3}{2} & \dfrac{3\pi}{2} & 0          & -1         & 0  & -\dfrac{3}{2} \\
    2            & 2\pi            & 1          & 0          & 2  & 0             \\ \hline
  \end{array}
  \]
  \end{ceqn}
\end{resolution}

%------------------------------------------------------------------------------%
\framedgraphic*[resolution]{Solução}{B_02_exemplo_5_curva_1}
\framedgraphic*[resolution]{Solução}{B_02_exemplo_5_curva_2}
\framedgraphic*[resolution]{Solução}{B_02_exemplo_5_curva_3}
\framedgraphic*[resolution]{Solução}{B_02_exemplo_5_curva_4}
\framedgraphic*[resolution]{Solução}{B_02_exemplo_5_curva_5}
\framedgraphic*[resolution]{Solução}{B_02_exemplo_5_curva_6}

%------------------------------------------------------------------------------%
